\section{Overlaps, runs, and deficits}

To understand the quantities in the lower bound, we first describe how
the cheapest overlaps fit together. The \emph{weight} of a step from a
permutation $u$ to a different permutation $v$ is the number of letters
that must be appended to a word ending in $u$ to make it end in $v$.
Thus the weight is $K$ minus the length of the longest suffix of $u$
that is a prefix of $v$.
Ordering one chosen occurrence of each permutation in a covering word
produces a Hamiltonian path whose total weight is at most the word's
length minus $K$. Conversely, appending the required letters along such
a path produces a covering word. Hence $S(K)$ is $K$ plus the minimum
weight of a Hamiltonian path in this overlap graph.

A step has weight one precisely when it passes to the next cyclic
rotation. We call the $K$ rotations of a permutation a \emph{rotation
class}. Traversing all of them in cyclic order gives a \emph{complete
block}. The class of $1234$, for instance, consists of $1234,2341,3412,4123$.
Within this block, each of the three steps costs one. Repeating the next
step would revisit $1234$.

Our arguments take place in paths of complete blocks whose rotation
classes are pairwise distinct. The minimum-path and forest
reduction produces auxiliary paths of this kind from an arbitrary
superpermutation~\cite{pantone-framework}. Its exact form is recalled
in Appendix~\ref{app:forest}. A superpermutation as written
need not itself traverse whole classes at once. We shall return to this
reduction after seeing what can be counted in one of its paths.

It is convenient to put $n=K-1$ and use the alphabet $\{0,\ldots,n\}$.
Suppose a block is entered at $(x_1,\ldots,x_n,s)$. After traversing
the block, the exit is $(s,x_1,\ldots,x_n)$. A weight-two step to a different rotation class must
enter at $(x_2,\ldots,x_n,x_1,s)$.
The other ordering of the two remaining symbols lies in the class just
visited. Thus the candidate in a different class rotates the first $n$ symbols
and leaves the last symbol fixed. It may still be excluded if that class
has appeared earlier in the path.
Successive weight-two steps therefore arrange the blocks themselves
around a cycle of length $n$.

Figure~\ref{fig:run} shows this at $K=5$. Each row is one block, and
the four rows are joined by weight-two steps. Continuing would bring us
back to the first row, so a class-distinct path can use at most four
of these blocks consecutively.

\begin{figure}[htbp]
\centering
\renewcommand{\arraystretch}{1.25}
\begin{tabular}{ccccc}
$12340$ & $23401$ & $34012$ & $40123$ & $01234$\\
$23410$ & $34102$ & $41023$ & $10234$ & $02341$\\
$34120$ & $41203$ & $12034$ & $20341$ & $03412$\\
$41230$ & $12304$ & $23041$ & $30412$ & $04123$
\end{tabular}
\caption{A full run on five symbols. Read each row from left to right,
then pass to the next row. Steps within a row have weight one; steps
between rows have weight two. The same weight-two step from the last
row returns to the first.}
\label{fig:run}
\end{figure}

A maximal succession of blocks joined at weight two is a \emph{run}.
If it contains $L$ blocks, its \emph{deficit} is $n-L$. A full run has
deficit zero. For example, stopping Figure~\ref{fig:run} after the
third row gives a run of deficit one. A path with $r$ runs and total
deficit $D$ thus contains $nr-D$ blocks.

This identity explains the importance of deficits. They measure how much
of the capacity of the runs goes unused. If the total number of blocks
is fixed, a larger deficit forces more runs, and passing from one run
to the next costs an extra letter.

For a precise illustration, suppose a single path visits all $T=n!$
rotation classes and every boundary between runs has weight three.
There are $T$ blocks, $T-r$ weight-two boundaries, and $r-1$ weight-three
boundaries. Its word length is therefore
\begin{equation}\label{eq:counting-example}
 K+(K-1)T+2(T-r)+3(r-1)=K+(K+1)T-3+r.
\end{equation}
Ignoring deficits would give at least $T/n=(n-1)!$ runs. Substituting
this into \eqref{eq:counting-example} gives $F_3(K)+K-3$. Thus the
familiar factorial terms already appear in this simple count; the extra
contribution comes from proving that runs leave some capacity unused.

The quantity we need to bound is the number of runs. Since
$T=nr-D$, an estimate $D\ge\delta r-\beta$ gives
$T\le(n-\delta)r+\beta$. In other words, apart from the endpoint
allowance $\beta$, each run covers at most $n-\delta$ classes on average.
The denominator $n-\delta$ in the eventual lower bound has this meaning.

The forest reduction extends this accounting to arbitrary words. It
handles paths that break into components and boundaries of weight greater
than three, while keeping the same consequence of a deficit estimate.
The estimate must hold on every consecutive segment of runs, because
the proof applies it after cutting paths and adding certain entrance
blocks. This is the scope of the local arguments below.

\section{The chains that resist the count}

For the remainder of the paper, assume $n\ge7$. We now ask how a path
can have many runs while keeping its deficit small.
Restrict attention to consecutive runs joined at weight three; call such
a path a \emph{genuine light segment}. Its blocks still belong to distinct
rotation classes. Record the segment by its sequence of run deficits.
For example, an entry $2$ represents a run missing two blocks, while an
entry $0$ represents a full run.

The local overlap restrictions prevent us from inserting arbitrarily
many zeros between small positive deficits. Put $a=n-2$. A string of
full runs has length at most $a$, unless the whole segment consists
of full runs, in which case its length is at most $a+1$. The pattern
$(0,1,0)$ is impossible: a run missing one block cannot have full runs
on both sides. More generally, a factor $(0,1,d)$ has $d\ge2$.
Finally, if a gap of $a$ full runs separates two positive deficits,
those adjacent deficits sum to at least $a+2$. These are the earlier
local run restrictions~\cite{pantone-framework}.

Thus the first small positive deficit that can sit alone between full
runs is $2$. Two such entries can be separated by at most $a-1$ zeros.
The pattern that is most troublesome for the deficit count consists of
many singleton twos separated by these long gaps. Each two contributes
the same deficit, while the zeros increase the number of runs being
covered by the estimate.

This leads to the chains used in the proof. In a factor starting and
ending in zero, a positive cluster is \emph{regular} if it is the
singleton $2$. Join consecutive regular clusters across gaps of $a-1$
zeros to form a maximal \emph{chain}. Its length is its number of
regular clusters. A gap of $a-2$ zeros joins consecutive chains into a
\emph{path of chains}; shorter gaps and other positive clusters cut the
path. At $a=5$, for example, the factor
\[
 0,\ \underbrace{2,0,0,0,0,2}_{\text{chain of length 2}},\
 \underbrace{0,0,0}_{\text{shortened gap}},\
 \underbrace{2}_{\text{chain of length 1}},\ 0
\]
has two chains. These definitions retain the original adjacencies:
cutting out a nonregular cluster does not join its two neighboring gaps.

There are two ways to improve the count. One can limit the length of
a single chain, or show that several long chains cannot follow one
another too efficiently. The earlier proof already does both. We shall
strengthen the second estimate. The form we need is
\begin{equation}\label{eq:CP}
 \ell_1+\cdots+\ell_q\le p+\frac t2(q-1)
\end{equation}
for every path of $q$ chains. Here $p$ is an endpoint allowance, while
$t/2$ is the limiting average chain length permitted by the estimate.
Write $\CP(p,t)$ if \eqref{eq:CP} holds in every zero-ended factor of
every genuine light segment.

Each boundary between chains has one fewer full run than a link within
a chain. For a fixed number of regular clusters, more such boundaries
therefore allow fewer zero entries without changing the total deficit.
This is why shortening the permitted average chain length strengthens
the deficit estimate. The earlier scalar argument makes this precise: all the nonregular clusters
and shorter gaps already supply enough deficit to cover their part of
the count. We recall that argument in Appendix~\ref{app:scalar}.

\subsection{Following a chain also relabels the symbols}

The deficit sequence alone does not record which classes have been
visited. To recover some of this information, follow a chain in actual
permutations. Denote the two cells by
\[
 C=(2,0^{a-1}),\qquad B=(2,0^{a-2}),
\]
where an exponent means repetition. A chain of length $\ell$ followed
by a shortened gap has pattern $T_\ell=C^{\ell-1}B$. The final $B$
consists of the chain's last deficit-two run and its outgoing shortened
gap; this description is used only when that gap is present.

The cells are particularly useful because their internal steps are
forced. To see what this means, relabel a current exit as
$(0,1,\ldots,n)$. A weight-three step leads to $(3,\ldots,n,i,j,k)$,
where $(i,j,k)$ is one of the six orders of $(0,1,2)$. The local cell
calculation forces the order $210$. After a completed $C$, the same
choice is forced if the next run has at least two blocks; after $B$,
it is forced if the next run has at least three. The alternatives
collide with classes in the cell just traversed~\cite{pantone-framework}.

Consequently, once the first entry is fixed, each cell carries us to a
specified final exit. Relabeling that exit back to $(0,1,\ldots,n)$
puts the next cell into the same coordinates. Call this relabeling
the cell's \emph{normalization}. It acts on the visited classes as well
as on the exit. This last point is essential: resetting the coordinates
does not erase the earlier part of the path.

In one-line notation on $0,\ldots,n$, the two normalizations are
\begin{equation}\label{eq:maps}
 g=(0,5,4,6,\ldots,n,2,3,1),\qquad
 \rho=(0,1,3,4,\ldots,n,2),\qquad
 N(C)=g,\quad N(B)=\rho g.
\end{equation}
Products act rightmost first, so the normalization of $T_\ell$ is
$h_\ell=\rho g^\ell$. The maps compose in reverse chronological
order. If a succession of completed chains has normalization equal to
the identity, its actual following short-run start repeats the first
start class. That is impossible in a genuine segment.

Let $p_0$ be the order of $g$. The earlier return calculation gives
\[
\begin{array}{c|c}
n&p_0\\\hline
3m&2m\\
3m+1&m(2m+1)\\
3m+2&(m+1)(2m+1).
\end{array}
\]
For example, at $n=7$ the nonzero cycles of $g$ are $(1,5,2,4,7)$
and $(3,6)$, giving order $10$. A single chain has length at most
$p_0$. If there is more than one chain in the path, the earlier
endpoint collisions reduce every chain's maximum to $b=p_0-1$
\cite{pantone-framework}.

Knowing these maxima does not yet give the best average. A very long
chain can force its successor to be short. The next two sections examine
the patterns that attain the resulting averages, beginning with a case
where there is just one efficient pair.

\subsection{The inherited restrictions on consecutive chains}

We collect the local geometric facts used in the counts below. The
parameters $p_0$ in this statement are the orders just displayed. These
are restrictions on actual occupied classes, not on arbitrary integer
sequences with the same deficit sums.

\begin{lemma}[Local chain restrictions~\cite{pantone-framework}]
\label{lem:inherited-chains}
Every chain has length at most $p_0$. In a path of at least two chains,
every length is at most $b=p_0-1$. The following additional restrictions
hold whenever the indicated consecutive chains occur.

For $n=3m+1$ and $m\ge3$, in two consecutive chains, if the first is longer than $m$, then the
second has length at most $m$. That successor has length at most $m-1$ if it is
not the last chain of the path.

For $n=3m+2$ and $m\ge2$, put $v=2m(m+1)$. In two consecutive chains, if the first is longer than $v$, then the
second has length at most $v-1$. If the first chain has length
at least $v+2$ and its successor is nonterminal, the successor has
length at most $v-2$.

For $n=3m$ and $m\ge3$, four consecutive chains cannot have their
middle two lengths equal to $b$ and both outer lengths at least two.
For $n=7$, the chain inequality $\CP(10,10)$ holds.
\end{lemma}

The cell forcing rules and these chain restrictions are the inherited
geometric inputs to the present argument. Their proofs identify a class
in common between the relevant cell paths. Further occupied classes can
exclude more paths, but cannot undo such a collision. This is why the
same restrictions apply inside every factor to which the forest count
will be applied.
